\section{Introduction}
\label{sec:intro}

Harvesting is among the most labour-intensive operations in fruit production, and the shrinking seasonal workforce has made robotic selective harvesting a long-standing goal of agricultural engineering \citep{bac2014harvesting,kootstra2021selective,zhou2022intelligent}. Systems for sweet pepper \citep{arad2020sweetpepper}, kiwifruit \citep{williams2019kiwifruit} and apple \citep{silwal2017apple} are pipelines of detection, localization, planning and grasping modules engineered for one crop, one decision rule and one scene layout. Vision--language--action (VLA) policies offer a different paradigm: a single network pretrained on web-scale vision--language data and large robot datasets maps images, proprioception and a natural-language instruction directly to actions \citep{brohan2023rt2,kim2024openvla,black2024pi0,pi05_2025}, so that \emph{which} fruits to pick, \emph{how many} and \emph{where} to put them become an instruction rather than an engineering effort.

\paragraph{Harvesting is a memory problem} A harvesting robot rarely has in view the information it needs for its next decision. Wind or the manipulator moves leaves over ripe fruits that were visible moments earlier \citep{hemming2014detectability}; grading, sampling and thinning require a given number of fruits; and the collection bin is usually outside the camera's view, so earlier picks cannot be re-observed. Forgetting has a direct agronomic cost: a forgotten ripe fruit is lost yield left to over-mature on the plant, a forgotten pick leads to over-harvesting that violates the quota, and a fruit whose ripeness was misjudged ends as unmarketable produce. In the benchmark introduced here, a leaf-hidden ripe fruit is invisible in a median 78\% of the frames in which it still hangs on the tree, so the decision to pick it must be taken largely from memory.

\paragraph{Time-indexed memory forgets more the longer the robot works} Memory-augmented VLA policies keep a bounded number of past frames, because every frame adds visual tokens and attention cost grows quadratically with their number \citep{vaswani2017attention}, and they fill this budget by sampling the history uniformly in time. This design has a structural flaw for manipulation. The frames that define a pick, the gripper closing on a fruit and the fruit detaching, span only about a dozen frames, and fewer than 5\% of uniformly sampled frames show the gripper within 3~cm of a fruit to be picked. With $B$ frames spread over $t$ elapsed frames, a pick is retained with probability roughly $(B-1)|S|/t$, where $|S|$ is the length of its key window, so coverage decays inversely with operating time (Section~\ref{sec:coverage}). At a 16-frame budget, uniform memory misses 43\% of completed picks, recovers the pick count at only 43\% of decisions and, after 600 frames, loses even the \emph{most recent} pick at 78\% of decisions. These losses matter: on a general manipulation-memory benchmark, removing the memory of the same base VLA reduces success from 66.7\% to 1.7\%, and halving its budget reduces success from 64.4\% to 26.7\% (Section~\ref{sec:res_general}).

\paragraph{Why the obvious remedies fall short} \emph{Enlarging the memory} treats the symptom: a uniform budget must grow linearly with operating time, at a quadratic attention cost that embedded field computers cannot afford. \emph{Better per-frame perception} cannot see what is hidden and mistakes a carried fruit for a new target; it recovers the correct pick count in only 62\% of our episodes. \emph{Learned memory selection or token compression} adds parameters and training signals, may behave differently in training and deployment, and is hard to audit. What is missing is a \emph{principle} that tells a harvesting robot what is worth remembering.

\paragraph{Key insight: index memory by events, not by time} In agricultural manipulation, task-relevant information is concentrated at a few physical interaction events. Every pick consists of an approach, a contact and a detachment; these moments change the state of the task, and they are physically observable, because the distance between the tool and the nearest fruit still to be picked drops below the fruit radius at contact and jumps at detachment. This mirrors human memory, which segments experience into discrete events and preferentially encodes their boundaries \citep{zacks2007event}, and whose episodic component binds \emph{what} happened to \emph{where} and \emph{when} \citep{tulving1972episodic,clayton1998episodic}.

\paragraph{HarvestMem} We introduce \emph{HarvestMem}, which equips a VLA harvesting policy with two complementary memories coupled through one physical signal (Fig.~\ref{fig:overview}). Its \textbf{spatial memory} answers \emph{where} the fruits are: a SAM-based fruit distance module (\fdm) segments fruits with the Segment Anything Model \citep{kirillov2023sam} while they are visible and anchors each one in 3D, persisting under foliage. Its \textbf{episodic memory} answers \emph{what} happened and \emph{when}: Picking-Event Keyframe Memory (\pekm) detects picks causally from the tool-to-fruit distance and reserves slots in the VLA's fixed memory for their contact and detachment keyframes. HarvestMem changes \emph{what} the policy remembers, not \emph{how much} or \emph{how}: \pekm\ adds no parameters, tokens or inference cost, contains the base model's uniform memory as the special case $k=0$, and is applied identically in training and deployment. Phase-aware recovery demonstrations (DART-H) complement the memory by addressing the compounding errors of closed-loop imitation \citep{ross2011dagger}.

\paragraph{Contributions} The main contributions of this work are:
\begin{enumerate}
\item \textbf{Problem formulation and benchmark.} We formulate selective harvesting with a VLA policy as a fixed-budget memory-allocation problem. We release \bench, a procedural ManiSkill3 orchard with a commercial 7-DoF RealMan RM75-B arm and two tasks that cannot be solved from the current view: picking leaf-hidden ripe fruits (FruitOccludedRipe) and meeting a quota with an out-of-view basket (FruitCount). The benchmark includes per-frame ground truth for events, occlusion and perception.
\item \textbf{Event-centric memory with a coverage analysis.} We identify a structural limitation of time-indexed memory: its coverage of interaction events decays inversely with operating time, and a simple sampling model predicts the measured coverage to within 0.012. We propose \pekm, which guarantees coverage of the most recent picks independently of episode length. It retains 99.6\% of completed picks against 57.0\% for uniform memory at the same budget, and with 8 frames it outperforms uniform memory with 16, at zero additional cost.
\item \textbf{Spatial--episodic coupling through promptable segmentation.} \fdm\ turns RGB-D images and proprioception into 3D fruit anchors, using visual (depth vacancy) and proprioceptive (grasp) evidence of picking, and reproduces the ground-truth picking events exactly. Its ablations establish that spatial memory acquired before occlusion is a \emph{precondition} for episodic memory: when anchors are built after occlusion, none of the 138 leaf-hidden ripe fruits is anchored and the occluded-task count accuracy falls from 1.00 to 0.00. A robustness analysis turns event memory into an explicit accuracy specification for perception modules.
\item \textbf{Phase-aware recovery demonstrations (DART-H).} Temporally correlated joint noise is injected only in free-motion phases, while clean planner actions are recorded as labels. This targets the covariate shift that limits closed-loop imitation in harvesting.
\item \textbf{A rigorous evaluation protocol.} Evaluation follows a protocol sealed before the formal runs. It uses paired bootstrap with Holm correction and equivalence margins, and audits at every decision that the frames fed to the model match the declared selection rule. Counterfactual controls separate the effect of event information from the mere perturbation of changing frames.
\end{enumerate}

The remainder of the paper is organized as follows. Section~\ref{sec:related} positions the work and Section~\ref{sec:benchmark} describes \bench. Section~\ref{sec:method} presents HarvestMem: the coverage analysis, \pekm, \fdm\ and DART-H. Section~\ref{sec:design} details the experimental design and Section~\ref{sec:results} reports the results. Section~\ref{sec:discussion} discusses implications and limitations for agricultural deployment, and Section~\ref{sec:conclusions} concludes.

\begin{figure}[t]
\centering
\fbox{\parbox{0.92\linewidth}{\centering\vspace{1cm}\tbd{Fig.~1 (to be drawn). Left: the memory needs of harvesting (ripe fruit hidden by wind-blown leaves; quota with an out-of-view basket); uniform memory misses the short pick events. Right: front/wrist RGB-D $\rightarrow$ \fdm\ (colour prompts, SAM masks, 3D anchors, vacancy + grasp evidence; \emph{where}) $\rightarrow$ $d_t$ $\rightarrow$ event detector $\rightarrow$ \pekm\ keyframes (\emph{what}/\emph{when}) $\rightarrow$ VLA memory expert $\rightarrow$ action chunk; the same selection function computes the memory of every training sample.}\vspace{1cm}}}
\caption{HarvestMem indexes memory by physical events instead of time. Fruits are anchored in 3D while visible (spatial memory), and the contact and detachment keyframes of picks are reserved in the VLA's fixed memory budget (episodic memory). The distance $d_t$ is the only interface between perception and the policy; it selects frame indices and never enters the network.}
\label{fig:overview}
\end{figure}
